% Generated by the blueprint tool (blueprint-pdf/1) from a blueprint snapshot; edit the
% Markdown sources, not this file.  Compile with `latexmk -pdf` or by running
% `pdflatex` twice (the second run settles the contents and the page references).
% Images, if any, are read from assets/ next to this file.
\documentclass[11pt]{article}
\usepackage[T1]{fontenc}
\usepackage[utf8]{inputenc}
\usepackage{lmodern}
\usepackage{textcomp}
\usepackage[a4paper,margin=2.5cm,headheight=14pt]{geometry}
\usepackage{amsmath,amssymb,amsthm}
\usepackage{amscd}  % KaTeX's CD diagrams; part of amsmath, so always there
\usepackage{mathtools}
\usepackage{bm}
\usepackage{graphicx}
\usepackage{xcolor}
\usepackage{array,booktabs,tabularx}
\usepackage{fancyvrb}
\usepackage[normalem]{ulem}
\usepackage{microtype}
% Optional, and only with their fonts: a small TeX installation has the
% packages without the fonts, which fails at the first formula.
\IfFileExists{mathrsfs.sty}{\IfFileExists{rsfs10.tfm}{\usepackage{mathrsfs}}{}}{}
\IfFileExists{stmaryrd.sty}{\IfFileExists{stmary10.tfm}{\usepackage{stmaryrd}}{}}{}
\IfFileExists{cancel.sty}{\usepackage{cancel}}{}
\usepackage[hyphens]{url}
\usepackage[pdfusetitle,bookmarksnumbered=false]{hyperref}
\usepackage{bookmark}
\definecolor{bplink}{HTML}{1F4E9C}
\definecolor{bpmuted}{HTML}{5F676E}
\hypersetup{colorlinks=true,linkcolor=bplink,urlcolor=bplink,citecolor=bplink}
\setlength{\parindent}{0pt}
\setlength{\parskip}{0.55em plus 0.2em minus 0.1em}
% Lists as tight as the page's.
\makeatletter
\def\@listi{\leftmargin\leftmargini\parsep 0pt\topsep 0.3em plus 0.1em\itemsep 0.1em plus 0.1em}
\let\@listI\@listi
\def\@listii{\leftmargin\leftmarginii\labelwidth\leftmarginii\advance\labelwidth-\labelsep\parsep 0pt\topsep 0.1em\itemsep 0.1em}
% An ordered list that does not start at 1.
\newcommand{\bpstartat}[1]{\setcounter{\@enumctr}{\numexpr#1-1\relax}}
\makeatother
\setcounter{tocdepth}{3}
\emergencystretch=3em
\allowdisplaybreaks

% A numbered statement, set as a paper sets one: "Theorem 2.1 (Title)."
% #1 is the lead ("Theorem~2.1"), #2 what follows it in roman.
\newenvironment{bpstatement}[2]{\par\addvspace{0.6\baselineskip}\noindent\textbf{#1}#2\textbf{.}\enspace\ignorespaces}
  {\par\addvspace{0.4\baselineskip}}
\newenvironment{bpstatementit}[2]{\begin{bpstatement}{#1}{#2}\itshape\ignorespaces}{\end{bpstatement}}
% A step: a relation with prose of its own, read under the object it starts from.
\newenvironment{bpstep}[1]{\par\addvspace{0.3\baselineskip}\begin{list}{}{\leftmargin=1.6em\rightmargin=0pt\topsep=0pt}\item\relax
  {\small\color{bpmuted}#1}\par\nopagebreak}{\end{list}\par}
% Notes about where an object is written out.
\newenvironment{bpnote}{\par\small\color{bpmuted}\itshape}{\par}
\newcommand{\bpomitted}[1]{\par{\small\color{bpmuted}\itshape #1\par}}
\newcommand{\bpbodyheading}[1]{\par\addvspace{0.4\baselineskip}\noindent\textbf{#1}\par\nopagebreak}
\newcommand{\bprule}{\par\medskip\noindent\hfil\rule{0.3\linewidth}{0.4pt}\hfil\par\medskip}
\newcommand{\bpboundary}[1]{\par\nopagebreak{\small\color{bpmuted}#1\par}\nopagebreak}
% A reference to something this export leaves out, and one that resolves nowhere.
\newcommand{\bpout}[1]{#1\textsuperscript{\hyperlink{bp-outside}{\ensuremath{\ast}}}}
\newcommand{\bpbroken}[1]{\textcolor{red!70!black}{[#1?]}}
\newcommand{\bpslug}[1]{\texttt{#1}}
% A bibliography entry: label, anchor, entry; the label hangs in the margin.
\newcommand{\bpbibitem}[3]{\par\setbox0\hbox{[#1]\enspace}\ifdim\wd0<5em\setbox0\hbox to 5em{[#1]\hfil}\fi
  \hangindent=\wd0 \hangafter=1 \noindent\box0\phantomsection\label{#2}#3\par\smallskip}
% A character the export knows no way to set.
\newcommand{\bpglyph}[1]{\fbox{\scriptsize #1}}
% Images, at their own size unless they are wider than the line.
\newsavebox{\bpimagebox}
\newcommand{\bpimage}[1]{\sbox{\bpimagebox}{\includegraphics{#1}}%
  \ifdim\wd\bpimagebox>\linewidth\resizebox{\linewidth}{!}{\usebox{\bpimagebox}}\else\usebox{\bpimagebox}\fi}
\newcommand{\bpmissingimage}[1]{\fbox{\parbox{0.8\linewidth}{\centering\small\color{bpmuted}Image not included: #1}}}
\DefineVerbatimEnvironment{bpverbatim}{Verbatim}{fontsize=\small,xleftmargin=1em}
\renewcommand{\qedsymbol}{\ensuremath{\blacksquare}}
\hypersetup{pdftitle={Weil Conjectures — C The Picard scheme of a curve},pdfsubject={Appendix C The Picard scheme of a curve; all levels},pdfcreator={blueprint}}
% KaTeX commands the website knows
\providecommand{\R}{\mathbb{R}}\providecommand{\N}{\mathbb{N}}\providecommand{\Z}{\mathbb{Z}}
\providecommand{\Q}{\mathbb{Q}}\providecommand{\C}{\mathbb{C}}
\providecommand{\Reals}{\mathbb{R}}\providecommand{\reals}{\mathbb{R}}\providecommand{\Complex}{\mathbb{C}}
\providecommand{\cnums}{\mathbb{C}}\providecommand{\natnums}{\mathbb{N}}\providecommand{\Bbb}[1]{\mathbb{#1}}
\providecommand{\bold}[1]{\mathbf{#1}}\providecommand{\frak}[1]{\mathfrak{#1}}\providecommand{\mathbbm}[1]{\mathbb{#1}}
\providecommand{\lang}{\langle}\providecommand{\rang}{\rangle}\providecommand{\bra}[1]{\langle #1\rvert}
\providecommand{\ket}[1]{\lvert #1\rangle}\providecommand{\braket}[1]{\langle #1\rangle}
\providecommand{\set}[1]{\{#1\}}\providecommand{\Set}[1]{\left\{#1\right\}}
\providecommand{\infin}{\infty}\providecommand{\isin}{\in}\providecommand{\plusmn}{\pm}\providecommand{\sub}{\subset}
\providecommand{\sube}{\subseteq}\providecommand{\supe}{\supseteq}\providecommand{\exist}{\exists}
\providecommand{\bull}{\bullet}\providecommand{\larr}{\leftarrow}\providecommand{\rarr}{\rightarrow}
\providecommand{\lrarr}{\leftrightarrow}\providecommand{\Larr}{\Leftarrow}\providecommand{\Rarr}{\Rightarrow}
\providecommand{\Lrarr}{\Leftrightarrow}\providecommand{\uarr}{\uparrow}\providecommand{\darr}{\downarrow}
\providecommand{\harr}{\leftrightarrow}\providecommand{\lArr}{\Leftarrow}\providecommand{\rArr}{\Rightarrow}
\providecommand{\hArr}{\Leftrightarrow}\providecommand{\alef}{\aleph}\providecommand{\alefsym}{\aleph}
\providecommand{\real}{\Re}\providecommand{\image}{\Im}\providecommand{\weierp}{\wp}
\providecommand{\Alpha}{\mathrm{A}}\providecommand{\Beta}{\mathrm{B}}\providecommand{\Epsilon}{\mathrm{E}}
\providecommand{\Zeta}{\mathrm{Z}}\providecommand{\Eta}{\mathrm{H}}\providecommand{\Iota}{\mathrm{I}}
\providecommand{\Kappa}{\mathrm{K}}\providecommand{\Mu}{\mathrm{M}}\providecommand{\Nu}{\mathrm{N}}
\providecommand{\Omicron}{\mathrm{O}}\providecommand{\Rho}{\mathrm{P}}\providecommand{\Tau}{\mathrm{T}}
\providecommand{\Chi}{\mathrm{X}}\providecommand{\omicron}{o}\providecommand{\gt}{>}\providecommand{\lt}{<}
\providecommand{\And}{\&}\providecommand{\hdashline}{\hline}\providecommand{\operatornamewithlimits}[1]{\operatorname*{#1}}
\providecommand{\llbracket}{[\![}\providecommand{\rrbracket}{]\!]}\providecommand{\cancel}[1]{#1}
\providecommand{\mathscr}{\mathcal}
% the characters beyond ASCII this document uses
\DeclareUnicodeCharacter{2160}{\textup{I}}
\DeclareUnicodeCharacter{2162}{\textup{III}}
% the project's macros ([katex.macros] in blueprint.toml)
\providecommand{\A}{}\renewcommand{\A}{\mathbf A}
\providecommand{\AEn}{}\renewcommand{\AEn}{A_{E,n}}
\providecommand{\Aaf}{}\renewcommand{\Aaf}{\mathbf A}
\providecommand{\Ab}{}\renewcommand{\Ab}{\mathsf{Ab}}
\providecommand{\AffSch}{}\renewcommand{\AffSch}{\operatorname{AffSch}}
\providecommand{\Aone}{}\renewcommand{\Aone}{\mathbf A^1}
\providecommand{\Bl}{}\renewcommand{\Bl}{\operatorname{Bl}}
\providecommand{\Br}{}\renewcommand{\Br}{\mathrm{Br}}
\providecommand{\C}{}\renewcommand{\C}{\mathbb C}
\providecommand{\Cech}{}\renewcommand{\Cech}{\check{H}}
\providecommand{\Coext}{}\renewcommand{\Coext}{\operatorname{Coext}}
\providecommand{\Cone}{}\renewcommand{\Cone}{\operatorname{Cone}}
\providecommand{\Cons}{}\renewcommand{\Cons}{\operatorname{Cons}}
\providecommand{\Cref}{}\renewcommand{\Cref}[1]{\text{#1}}
\providecommand{\D}{}\renewcommand{\D}{D}
\providecommand{\DLAFT}{}\renewcommand{\DLAFT}{\operatorname{FT}}
\providecommand{\DLAFr}{}\renewcommand{\DLAFr}{\operatorname{Fr}}
\providecommand{\DLAQ}{}\renewcommand{\DLAQ}{{\overline{\mathbf Q}_\ell}}
\providecommand{\DLRQ}{}\renewcommand{\DLRQ}{{\overline{\mathbf Q}_\ell}}
\providecommand{\Dbc}{}\renewcommand{\Dbc}{D^b_c}
\providecommand{\Dc}{}\renewcommand{\Dc}{D_{\mathrm{cons}}}
\providecommand{\Dctf}{}\renewcommand{\Dctf}{D^b_{ctf}}
\providecommand{\Dic}{}\renewcommand{\Dic}{D_{\IndCons}}
\providecommand{\Dplus}{}\renewcommand{\Dplus}{D^{+}}
\providecommand{\End}{}\renewcommand{\End}{\operatorname{End}}
\providecommand{\Et}{}\renewcommand{\Et}{\mathrm{\acute{e}t}}
\providecommand{\Ex}{}\renewcommand{\Ex}{\mathrm{Ex}}
\providecommand{\Ext}{}\renewcommand{\Ext}{\operatorname{Ext}}
\providecommand{\F}{}\renewcommand{\F}{\mathbb F}
\providecommand{\FT}{}\renewcommand{\FT}{\operatorname{FT}}
\providecommand{\Fbar}{}\renewcommand{\Fbar}{{\overline{\mathbf F}_q}}
\providecommand{\Fell}{}\renewcommand{\Fell}{\mathbf F_\ell}
\providecommand{\Fib}{}\renewcommand{\Fib}{\operatorname{Fib}}
\providecommand{\Fix}{}\renewcommand{\Fix}{\operatorname{Fix}}
\providecommand{\Fp}{}\renewcommand{\Fp}{{\mathbf F_{p}}}
\providecommand{\Fq}{}\renewcommand{\Fq}{{\mathbf F_q}}
\providecommand{\Fqn}{}\renewcommand{\Fqn}{{\mathbf F_{q^{n}}}}
\providecommand{\Frob}{}\renewcommand{\Frob}{\operatorname{Frob}}
\providecommand{\GL}{}\renewcommand{\GL}{\operatorname{GL}}
\providecommand{\Gal}{}\renewcommand{\Gal}{\operatorname{Gal}}
\providecommand{\Hens}{}\renewcommand{\Hens}{\operatorname{Hens}}
\providecommand{\Hilb}{}\renewcommand{\Hilb}{\operatorname{Hilb}}
\providecommand{\Hom}{}\renewcommand{\Hom}{\operatorname{Hom}}
\providecommand{\IndCons}{}\renewcommand{\IndCons}{\operatorname{IndCons}}
\providecommand{\Kar}{}\renewcommand{\Kar}{\operatorname{Kar}}
\providecommand{\LT}{}\renewcommand{\LT}{\operatorname{LT}}
\providecommand{\Lie}{}\renewcommand{\Lie}{\operatorname{Lie}}
\providecommand{\Loc}{}\renewcommand{\Loc}{\operatorname{Loc}}
\providecommand{\Mod}{}\renewcommand{\Mod}{\mathsf{Mod}}
\providecommand{\N}{}\renewcommand{\N}{\mathbb{N}}
\providecommand{\Needspace}{}\renewcommand{\Needspace}{\par}
\providecommand{\O}{}\renewcommand{\O}{{\mathcal{O}}}
\providecommand{\Ob}{}\renewcommand{\Ob}{\mathrm{Ob}}
\providecommand{\Oe}{}\renewcommand{\Oe}{\mathcal O_E}
\providecommand{\PP}{}\renewcommand{\PP}{\mathbf{P}}
\providecommand{\Perf}{}\renewcommand{\Perf}{\operatorname{Perf}}
\providecommand{\Pic}{}\renewcommand{\Pic}{\mathrm{Pic}}
\providecommand{\Pm}{}\renewcommand{\Pm}{\mathbf P^{m}}
\providecommand{\Pone}{}\renewcommand{\Pone}{\mathbf P^{1}}
\providecommand{\Pqn}{}\renewcommand{\Pqn}{\mathbf P^{n}}
\providecommand{\ProFin}{}\renewcommand{\ProFin}{\operatorname{ProFin}}
\providecommand{\Proet}{}\renewcommand{\Proet}{\mathrm{pro\acute{e}t}}
\providecommand{\Q}{}\renewcommand{\Q}{\mathbb Q}
\providecommand{\Qb}{}\renewcommand{\Qb}{\overline{\mathbf Q}_{\ell}}
\providecommand{\Qbar}{}\renewcommand{\Qbar}{\overline{\mathbf Q}}
\providecommand{\QbarX}{}\renewcommand{\QbarX}{\overline{\mathbf Q}_{\ell,X}}
\providecommand{\Qbarl}{}\renewcommand{\Qbarl}{{\overline{\mathbf Q}_{\ell}}}
\providecommand{\Qell}{}\renewcommand{\Qell}{{\mathbf{Q}_{\ell}}}
\providecommand{\Ql}{}\renewcommand{\Ql}{\mathbf Q_{\ell}}
\providecommand{\Qlb}{}\renewcommand{\Qlb}{{\overline{\mathbf Q}_\ell}}
\providecommand{\Qlbar}{}\renewcommand{\Qlbar}{\overline{\mathbf Q}_\ell}
\providecommand{\R}{}\renewcommand{\R}{\mathbb{R}}
\providecommand{\RG}{}\renewcommand{\RG}{\operatorname{\mathrm{R}\Gamma}}
\providecommand{\RGamma}{}\renewcommand{\RGamma}{R\Gamma}
\providecommand{\RGammac}{}\renewcommand{\RGammac}{R\Gamma_c}
\providecommand{\RHom}{}\renewcommand{\RHom}{R\mathcal{H}om}
\providecommand{\RHomX}{}\renewcommand{\RHomX}{R\mathcal H om_X}
\providecommand{\Rat}{}\renewcommand{\Rat}{\operatorname{IsRat}}
\providecommand{\Res}{}\renewcommand{\Res}{\operatorname{Res}}
\providecommand{\Rlim}{}\renewcommand{\Rlim}{R\!\varprojlim}
\providecommand{\SIm}{}\renewcommand{\SIm}{\operatorname{SIm}}
\providecommand{\Sch}{}\renewcommand{\Sch}{\mathsf{Sch}}
\providecommand{\Sh}{}\renewcommand{\Sh}{\operatorname{Sh}}
\providecommand{\Span}{}\renewcommand{\Span}{\operatorname{span}}
\providecommand{\Spec}{}\renewcommand{\Spec}{\operatorname{Spec}}
\providecommand{\Tr}{}\renewcommand{\Tr}{\operatorname{Tr}}
\providecommand{\Trsm}{}\renewcommand{\Trsm}{\operatorname{Tr}^{\mathrm{sm}}}
\providecommand{\Wt}{}\renewcommand{\Wt}{\operatorname{Wt}}
\providecommand{\Z}{}\renewcommand{\Z}{\mathbb{Z}}
\providecommand{\Zar}{}\renewcommand{\Zar}{\mathrm{zar}}
\providecommand{\Zell}{}\renewcommand{\Zell}{{\mathbf{Z}_{\ell}}}
\providecommand{\Zl}{}\renewcommand{\Zl}{{\mathbf Z_\ell}}
\providecommand{\affet}{}\renewcommand{\affet}[1]{{#1}^{\mathrm{aff}}_{\mathrm{ét}}}
\providecommand{\affproet}{}\renewcommand{\affproet}[1]{{#1}^{\mathrm{aff}}_{\mathrm{proét}}}
\providecommand{\boxt}{}\renewcommand{\boxt}{\boxtimes_A^{\mathrm L}}
\providecommand{\cF}{}\renewcommand{\cF}{\mathcal{F}}
\providecommand{\cHom}{}\renewcommand{\cHom}{R\mathcal{H}\!om}
\providecommand{\cI}{}\renewcommand{\cI}{\mathcal{I}}
\providecommand{\cL}{}\renewcommand{\cL}{\mathcal L}
\providecommand{\cM}{}\renewcommand{\cM}{\mathcal M}
\providecommand{\cN}{}\renewcommand{\cN}{\mathcal N}
\providecommand{\cO}{}\renewcommand{\cO}{\mathcal O}
\providecommand{\cT}{}\renewcommand{\cT}{\mathcal T}
\providecommand{\checked}{}\renewcommand{\checked}{\color{red}\aftergroup\restorecolor}
\providecommand{\coeff}{}\renewcommand{\coeff}{\operatorname{coeff}}
\providecommand{\coker}{}\renewcommand{\coker}{\operatorname{coker}}
\providecommand{\colim}{}\renewcommand{\colim}{\operatorname{colim}}
\providecommand{\constructedby}{}\renewcommand{\constructedby}[1]{\par\smallskip\noindent\textcolor{purple!60!black}{\textsf{Constructed by:} }
  \footnotesize\sloppy\nolinkurl{#1}\par}
\providecommand{\deepresult}{}\renewcommand{\deepresult}[1]{\par\smallskip\noindent\textcolor{orange!45!black}{
  \textsf{Blueprint role: named deep theorem node; see the cluster proof architecture.}}
  \par\noindent\textcolor{orange!45!black}{\textsf{Mathematical source:} }#1\par}
\providecommand{\dt}{}\renewcommand{\dt}{\mathbin{\otimes^{\mathbf L}}}
\providecommand{\dw}{}\renewcommand{\dw}{\operatorname{dw}}
\providecommand{\et}{}\renewcommand{\et}[1]{{#1}_{\mathrm{ét}}}
\providecommand{\ev}{}\renewcommand{\ev}{\operatorname{ev}}
\providecommand{\fib}{}\renewcommand{\fib}{\operatorname{fib}}
\providecommand{\formalizationnote}{}\renewcommand{\formalizationnote}[1]{\begin{quote}\small\raggedright\textbf{Formalization note.} #1\end{quote}}
\providecommand{\fpqc}{}\renewcommand{\fpqc}{X_{\mathrm{fpqc}}}
\providecommand{\gammacont}{}\renewcommand{\gammacont}{\Gamma_{\mathrm{cont}}}
\providecommand{\gm}{}\renewcommand{\gm}{\mathbf G_m}
\providecommand{\hocolim}{}\renewcommand{\hocolim}{\operatorname{hocolim}}
\providecommand{\holim}{}\renewcommand{\holim}{\operatorname{holim}}
\providecommand{\hrB}{}\renewcommand{\hrB}{\@ifnextchar]{\hrEnd}{\hrStep}}
\providecommand{\hrEnd}{}\renewcommand{\hrEnd}{\@firstoftwo{}}
\providecommand{\hrStep}{}\renewcommand{\hrStep}{\@firstoftwo{\hrB}}
\providecommand{\id}{}\renewcommand{\id}{\operatorname{id}}
\providecommand{\kunFl}{}\renewcommand{\kunFl}{\mathbf F_\ell}
\providecommand{\kunZn}{}\renewcommand{\kunZn}{\mathbf Z/n\mathbf Z}
\providecommand{\kuncd}{}\renewcommand{\kuncd}{\operatorname{cd}}
\providecommand{\kundim}{}\renewcommand{\kundim}{\dim}
\providecommand{\kunotimes}{}\renewcommand{\kunotimes}[1]{\otimes^{\mathbf L}_{#1}}
\providecommand{\kunsupp}{}\renewcommand{\kunsupp}{\operatorname{Supp}}
\providecommand{\lotimes}{}\renewcommand{\lotimes}{\otimes^{\mathrm L}}
\providecommand{\m}{}\renewcommand{\m}{{\mathfrak{m}}}
\providecommand{\n}{}\renewcommand{\n}{{\mathfrak{n}}}
\providecommand{\nagataGap}{}\renewcommand{\nagataGap}[1]{\par\noindent\textbf{Formalization gap.} #1\par}
\providecommand{\openobligation}{}\renewcommand{\openobligation}[1]{\par\smallskip\noindent\textcolor{red!65!black}{\textsf{Open obligation:} }
  \footnotesize\sloppy\nolinkurl{#1}\par}
\providecommand{\p}{}\renewcommand{\p}{{\mathfrak{p}}}
\providecommand{\pr}{}\renewcommand{\pr}{\operatorname{pr}}
\providecommand{\preshv}{}\renewcommand{\preshv}[1]{\mathrm{PreShv}(#1)}
\providecommand{\pro}{}\renewcommand{\pro}[1]{\mathrm{Pro}(#1)}
\providecommand{\proet}{}\renewcommand{\proet}{X_{\mathrm{proét}}}
\providecommand{\q}{}\renewcommand{\q}{{\mathfrak{q}}}
\providecommand{\rad}{}\renewcommand{\rad}{\operatorname{rad}}
\providecommand{\rank}{}\renewcommand{\rank}{\operatorname{rank}}
\providecommand{\red}{}\renewcommand{\red}{\color{red}}
\providecommand{\restorecolor}{}\renewcommand{\restorecolor}{\color{black}}
\providecommand{\retiredobligation}{}\renewcommand{\retiredobligation}[1]{\par\smallskip\noindent\textcolor{red!75!black}{\textsf{Retired invalid route:} }
  \footnotesize\sloppy\nolinkurl{#1}\par}
\providecommand{\rk}{}\renewcommand{\rk}{\operatorname{rk}}
\providecommand{\sep}{}\renewcommand{\sep}{\textnormal{sep}}
\providecommand{\shv}{}\renewcommand{\shv}[1]{\mathrm{Shv}(#1)}
\providecommand{\spec}{}\renewcommand{\spec}[1]{\mathrm{Spec}(#1)}
\providecommand{\stacks}{}\renewcommand{\stacks}[1]{\href{https://stacks.math.columbia.edu/tag/#1}{\textnormal{Stacks, Tag~#1}}}
\providecommand{\stacksproject}{}\renewcommand{\stacksproject}[1]{{\cite[\href{https://stacks.math.columbia.edu/tag/#1}{Tag #1}]{stacks-project}}}
\providecommand{\statusnote}{}\renewcommand{\statusnote}[1]{\begin{quote}\small\raggedright\textbf{Status.} #1\end{quote}}
\providecommand{\tildeideal}{}\renewcommand{\tildeideal}[2]{\widetilde{(#1, #2)}}
\providecommand{\tildering}{}\renewcommand{\tildering}[3]{{#1}_{\widetilde{(#2, #3)}}}
\providecommand{\tr}{}\renewcommand{\tr}{\operatorname{tr}}
\providecommand{\wt}{}\renewcommand{\wt}{\operatorname{wt}}

\begin{document}
\pagestyle{plain}
\begin{center}
{\LARGE\bfseries Weil Conjectures\par}
\vspace{0.7em}
{\large Appendix C: The Picard scheme of a curve\par}
\vspace{0.5em}
{\small\color{bpmuted}Snapshot 2026-09-26, revision 8e4cb88e2a\par}
\end{center}
\vspace{0.5em}
{\small\tableofcontents}
\bigskip\hrule\bigskip

% object: chap-The-Picard-scheme-of-a-curve (blueprint/chap-The-Picard-scheme-of-a-curve/_section.md)
\section*{\texorpdfstring{C\quad The Picard scheme of a curve}{C The Picard scheme of a curve}}
\phantomsection\label{bp1}%
\addcontentsline{toc}{section}{\texorpdfstring{C\quad The Picard scheme of a curve}{C The Picard scheme of a curve}}
% text: prose-C-01
\phantomsection\label{bp2}%
For a curve \(C\) proper and connected over a separably closed field and \(n\) invertible in that field, the Kummer sequence shows that \(H^1(C,\mu_n)\) is the \(n\)-torsion of the Picard group of \(C\). But the Picard group is an abstract group. In particular, nothing about it says that its \(n\)-torsion is finite, let alone isomorphic to \((\mathbf Z/n)^{2g}\). Nor does it say why Frobenius should have no fixed vector on the \(\ell\)-adic limit, as Proposition \bpout{14.3} requires. Both questions become statements of Appendix \bpout{B} once we know that \(\operatorname{Pic}^0(C)\), the group of invertible sheaves of degree zero, is the group of points of an abelian variety.

Over \(\mathbf C\) this is the Abel–Jacobi theorem, \(\operatorname{Pic}^0(C)\simeq
H^0(C,\Omega^1)^\vee/H_1(C,\mathbf Z)\). Over an arbitrary field we construct the Jacobian as a Picard scheme, which represents invertible sheaves in families.

Let \(C\) be a smooth proper connected curve of genus \(g\) over a separably closed field \(k\). The idea is that most invertible sheaves of degree \(g\) are given by divisors. Suppose \(\mathcal L\) has degree \(g\) and \(h^0(\mathcal L)=1\). Then the zeros of its unique section up to scalars form an effective divisor of degree \(g\), and this divisor determines \(\mathcal L\). Effective divisors of degree \(g\) are the points of the Hilbert scheme \(\operatorname{Hilb}^{g}_{C/k}\), a smooth proper variety of dimension \(g\) (Theorem \hyperref[bp21]{C.6}).

By Riemann–Roch, every invertible sheaf of degree \(g\) has \(h^0\geq1\), with equality exactly when \(h^1=0\). This condition is open in families, by semicontinuity, and it defines an open chart \(W\subset\operatorname{Hilb}^{g}_{C/k}\) of the Picard functor (Proposition \hyperref[bp24]{C.7}). Riemann–Roch also moves every invertible sheaf into this chart by a twist defined over \(k\) (Lemma \hyperref[bp27]{C.8}). Hence the translates of \(W\) cover the Picard functor, and they glue to the Picard scheme. Its components are proper because the Abel map \(\operatorname{Hilb}^{g}_{C/k}\to\operatorname{Pic}^g\) is surjective and its source is proper.

The Picard functor itself is a sheafified quotient of \(T\mapsto\operatorname{Pic}(C\times T)\). We have to divide out the invertible sheaves pulled back from \(T\), since these are trivial locally on \(T\). When \(C\) has a rational point, we describe the Picard functor by invertible sheaves trivialized along that point (Lemma \hyperref[bp7]{C.2}).

The outcome is Theorem \hyperref[bp36]{C.11}: the component of degree zero, the Jacobian \(J(C)\), is an abelian variety of dimension \(g\) with \(J(C)(k)=\operatorname{Pic}^0(C)\) and tangent space \(H^1(C,\mathcal O_C)\) at the origin. Over a finite field we descend the Jacobian to an abelian variety \(J_0\) over \(\mathbf F_q\). Then Propositions \hyperref[bp49]{C.16} and \hyperref[bp52]{Proposition C.17} give an isomorphism \(\varprojlim_nH^1(C,\mathbf Z/\ell^n(1))\simeq T_\ell J_0\) compatible with Frobenius, which is the form in which Proposition \bpout{14.3} uses the Jacobian; see also [\hyperref[bp54]{Mil86b}].

% object: sec-The-Picard-functor (blueprint/chap-The-Picard-scheme-of-a-curve/sec-The-Picard-functor/_section.md), content hidden
\subsection*{\texorpdfstring{C.1\quad The Picard functor}{C.1 The Picard functor}}
\phantomsection\label{bp3}%
\addcontentsline{toc}{subsection}{\texorpdfstring{C.1\quad The Picard functor}{C.1 The Picard functor}}
% text: prose-C-02
\phantomsection\label{bp4}%
We want a scheme whose points over \(T\) are families of invertible sheaves on \(C\) parametrized by \(T\). The obvious candidate \(T\mapsto\operatorname{Pic}(C_T)\), however, contains too much. For example, for \(C=\mathbf P^1_k\) and connected \(T\), every invertible sheaf on \(\mathbf P^1_T\) is \(\mathcal O(d)\otimes f_T^*\mathcal M\) for a unique \(d\in\mathbf Z\) and a unique \(\mathcal M\in\operatorname{Pic}(T)\). In other words \(\operatorname{Pic}(\mathbf P^1_T)=\mathbf Z\oplus\operatorname{Pic}(T)\), whereas the Picard scheme should be the constant scheme \(\mathbf Z\). In fact the functor \(T\mapsto\operatorname{Pic}(C_T)\) is not even a Zariski sheaf, because the pullback \(f_T^*\mathcal M\) of a nontrivial invertible sheaf \(\mathcal M\) on \(T\) is trivial locally on \(T\) but not globally, so no scheme represents this functor.

The Picard functor is therefore defined as the sheafification of \(T\mapsto\operatorname{Pic}(C_T)/f_T^*\operatorname{Pic}(T)\) (Definition \hyperref[bp5]{C.1}). When \(C\) has a section, Lemma \hyperref[bp7]{C.2} removes the pullbacks from \(T\) in another way, by fixing a trivialization along the section. We use this rigidified description to construct the chart of Proposition \hyperref[bp24]{C.7} and the Picard scheme of Theorem \hyperref[bp36]{C.11}.

% object: definition-relative-picard-functor (blueprint/chap-The-Picard-scheme-of-a-curve/sec-The-Picard-functor/definition-relative-picard-functor/_section.md)
\phantomsection\label{bp5}%
\begin{bpstatement}{Definition~C.1}{ (Relative Picard functor)}
For a morphism of schemes \(f:C\to S\), let

\[
\operatorname{Pic}_{C/S}
=\bigl(T\longmapsto
\operatorname{Pic}(C_T)/f_T^*\operatorname{Pic}(T)\bigr)^{\mathrm{fppf}}.
\]

Here \(C_T=C\times_ST\), and the superscript denotes sheafification on the big fppf site of \(S\). When this functor is represented by a scheme, we call the representing scheme the \emph{Picard scheme}.
\end{bpstatement}

% text: prose-C-03
\phantomsection\label{bp6}%
Equivalently, one may sheafify \(T\mapsto\operatorname{Pic}(C_T)\), since an invertible sheaf on \(T\) becomes trivial on a Zariski cover. The sheafification is needed: invertible sheaves satisfy effective descent, but their automorphisms can prevent their isomorphism classes from forming a sheaf. For this reason neither the displayed quotient nor \(T\mapsto\operatorname{Pic}(C_T)\) need be an fppf sheaf [Stacks, \href{https://stacks.math.columbia.edu/tag/0B9L}{0B9L}]. If \(f\) has a section, the next lemma describes the sheafification directly. We also write \(\operatorname{Pic}_{C/S}\) for the Picard scheme when it exists.

% object: lemma-rigidification-along-section (blueprint/chap-The-Picard-scheme-of-a-curve/sec-The-Picard-functor/lemma-rigidification-along-section/_section.md)
\phantomsection\label{bp7}%
\begin{bpstatementit}{Lemma~C.2}{ (Rigidification along a section)}
Suppose that \(f\) has a section \(\sigma:S\to C\) and that

\[
\mathcal O_T\longrightarrow f_{T,*}\mathcal O_{C_T}
\quad\text{is an isomorphism for every }T\to S.
\tag{$\ast$}
\]

Then \(\operatorname{Pic}_{C/S}(T)\) is identified, functorially in \(T\), with the set of isomorphism classes of pairs \((\mathcal L,\alpha)\), with \(\mathcal L\) invertible on \(C_T\) and \(\alpha:\sigma_T^*\mathcal L\simeq\mathcal O_T\). Every such pair has only the identity automorphism. There is a split exact sequence

\[
0\longrightarrow\operatorname{Pic}(T)
\xrightarrow{f_T^*}\operatorname{Pic}(C_T)
\longrightarrow\operatorname{Pic}_{C/S}(T)
\longrightarrow0,
\]

whose retraction is \(\sigma_T^*\). The normalization map is

\[
\mathcal L\longmapsto
\mathcal L\otimes f_T^*(\sigma_T^*\mathcal L)^{-1}.
\]
\end{bpstatementit}

% proof: proof/lemma-rigidification-along-section
\begin{proof}[{Proof}]\phantomsection\label{bp8}%
The point is that rigidified pairs have no automorphisms other than the identity. Because of this they form a sheaf, and the quotient presheaf maps isomorphically onto this sheaf.

An automorphism of an invertible sheaf is multiplication by a global unit. By condition \((\ast)\) this unit comes from a unit on \(T\), and compatibility with a rigidification forces it to be \(1\). To see that the rigidified functor is a sheaf, suppose we are given rigidified classes on an fppf cover that agree on overlaps. We then choose local isomorphisms and rescale them so that they preserve the rigidifications. Such isomorphisms are unique, since rigidified pairs have no nontrivial automorphisms. In particular they satisfy the cocycle condition, and the pairs glue by effective descent for invertible sheaves.

The normalization map identifies the quotient presheaf \(T\mapsto\operatorname{Pic}(C_T)/f_T^*\operatorname{Pic}(T)\) with this rigidified functor. Since the latter is already a sheaf, the quotient presheaf needs no sheafification and equals \(\operatorname{Pic}_{C/S}\). The sequence splits since \(\sigma_T^*f_T^*=\mathrm{id}\); for this argument see [Stacks, \href{https://stacks.math.columbia.edu/tag/0B9N}{0B9N}].
\end{proof}

% text: prose-C-04
\phantomsection\label{bp9}%
For a smooth proper geometrically connected curve \(C\) over a field \(k\), we have \(H^0(C,\mathcal O_C)=k\). For every field extension \(k'/k\) the curve \(C_{k'}\) is again smooth, proper and integral, and flat base change for coherent cohomology gives \(H^0(C_{k'},\mathcal O_{C_{k'}})=k'\) and \(\dim_{k'}H^1(C_{k'},\mathcal O_{C_{k'}})=\dim_kH^1(C,\mathcal O_C)\). Every \(k\)-scheme is flat over \(k\), so the same flat base change shows that \((\ast)\) holds for \(C\to\operatorname{Spec}k\). If \(C\) has a \(k\)-point \(P\), the preceding lemma applies with \(\sigma=P\) (see [Stacks, \href{https://stacks.math.columbia.edu/tag/0B9U}{0B9U}]).


% object: sec-Riemann-Roch-and-Serre-duality (blueprint/chap-The-Picard-scheme-of-a-curve/sec-Riemann-Roch-and-Serre-duality/_section.md), content hidden
\subsection*{\texorpdfstring{C.2\quad Riemann–Roch and Serre duality}{C.2 Riemann–Roch and Serre duality}}
\phantomsection\label{bp10}%
\addcontentsline{toc}{subsection}{\texorpdfstring{C.2\quad Riemann–Roch and Serre duality}{C.2 Riemann–Roch and Serre duality}}
% text: prose-C-05
\phantomsection\label{bp11}%
The chart of the Picard scheme in degree \(g\) is cut out by the condition \(h^1=0\). Let \(\mathcal L\) be an invertible sheaf of degree \(g\) on a curve of genus \(g\). Then Riemann–Roch gives \(h^0(\mathcal L)-h^1(\mathcal L)=1\), so \(h^1(\mathcal L)=0\) if and only if \(\mathcal L\) has exactly one section up to scalars. In a family over \(T\), coherent base change shows that the condition “\(Rf_{T,*}\mathcal L\) is an invertible sheaf in degree zero” is open on \(T\) and compatible with base change. It therefore defines an open subfunctor of the Picard functor.

The condition is not closed, because \(h^0\) can jump up under specialization. For example, on a curve of genus two a divisor \(D\) of degree two has \(h^0(\mathcal O(D))=1\) unless \(\mathcal O(D)\simeq\omega_{C/k}\), in which case \(h^0=2\). The chart therefore misses the line of canonical divisors in \(\operatorname{Hilb}^{2}_{C/k}\), and we need its translates to cover the Picard scheme in degree two.

Fix a smooth proper geometrically connected curve \(C\) over \(k\), with structure morphism \(f:C\to\operatorname{Spec}k\), and set \(g=\dim_kH^1(C,\mathcal O_C)\). Such a curve is projective [Stacks, \href{https://stacks.math.columbia.edu/tag/0A26}{0A26}].

% object: theorem-riemann-roch-coherent-serre-duality (blueprint/chap-The-Picard-scheme-of-a-curve/sec-Riemann-Roch-and-Serre-duality/theorem-riemann-roch-coherent-serre-duality/_section.md)
\phantomsection\label{bp12}%
\begin{bpstatementit}{Theorem~C.3}{ (Riemann–Roch and coherent Serre duality)}
For an invertible sheaf \(\mathcal L\) on \(C\) we have

\[
\chi(C,\mathcal L)=\deg\mathcal L+1-g,
\qquad
H^1(C,\mathcal L)^\vee
\simeq H^0(C,\omega_{C/k}\otimes\mathcal L^{-1}),
\qquad \deg\omega_{C/k}=2g-2,
\qquad \dim_kH^0(C,\omega_{C/k})=g.
\]

In particular, if \(\deg\mathcal L>2g-2\), then

\[
H^1(C,\mathcal L)=0,
\qquad h^0(C,\mathcal L)=\deg\mathcal L+1-g.
\]
\end{bpstatementit}

% text: prose-C-06
\phantomsection\label{bp13}%
For \(C=\mathbf P^1_k\) and \(\mathcal L=\mathcal O(d)\) with \(d\geq0\), the sections are the binary forms of degree \(d\), a space of dimension \(d+1=\deg\mathcal L+1-0\), and \(\omega=\mathcal O(-2)\) has degree \(-2=2g-2\).

We quote the theorem from [Stacks, \href{https://stacks.math.columbia.edu/tag/0B5B}{0B5B}], where it is proved for smooth curves and the duality is given by the residue trace on differentials. We prove its compatibility with field extension in Corollary \hyperref[bp14]{C.4} below. For the displayed vanishing, recall that a nonzero section of an invertible sheaf on an integral curve has an effective divisor of zeros. Hence an invertible sheaf of negative degree has no nonzero section. If \(\deg\mathcal L>2g-2\), then \(\omega_{C/k}\otimes\mathcal L^{-1}\) has negative degree, and \(H^1(C,\mathcal L)=0\) by duality. The dimension formula then follows from Riemann–Roch. For \(\mathcal L=\mathcal O_C\) the duality gives \(\dim_kH^0(C,\omega_{C/k})=\dim_kH^1(C,\mathcal O_C)=g\).

The Picard scheme is constructed over a separably closed field, but its points are tested over arbitrary extensions. Hence Riemann–Roch and coherent base change must be compatible with extension of the ground field, and for duality we need this for the pairing itself, not only for the dimensions.

% object: corollary-serre-duality-field-extension (blueprint/chap-The-Picard-scheme-of-a-curve/sec-Riemann-Roch-and-Serre-duality/corollary-serre-duality-field-extension/_section.md)
\phantomsection\label{bp14}%
\begin{bpstatementit}{Corollary~C.4}{ (Serre duality and field extension)}
Let \(k'/k\) be a field extension. Under flat base change for coherent cohomology and the identification \(\omega_{C_{k'}/k'}\simeq\omega_{C/k}|_{C_{k'}}\), the Serre duality pairing for \(\mathcal L|_{C_{k'}}\) is obtained from that for \(\mathcal L\) by extension of scalars. In particular, the genus is unchanged by field extension.
\end{bpstatementit}

% proof: proof/corollary-serre-duality-field-extension
\begin{proof}[{Proof}]\phantomsection\label{bp15}%
We factor the duality on \(C\) through a finite map to \(\mathbf P^1_k\), where the trace is an explicit residue, and check that each step is compatible with extension of scalars.

A nonconstant rational function on \(C\) extends to a finite map \(\phi:C\to\mathbf P^1_k\), which is flat because \(\phi_*\mathcal O_C\) is torsion free over the discrete valuation rings of \(\mathbf P^1_k\). We may construct the duality on \(C\) as duality for \(\phi\), given by the dual of the finite locally free module \(\phi_*\mathcal L\), followed by duality on \(\mathbf P^1_k\), whose trace \(H^1(\mathbf P^1_k,\omega)\to k\) takes the coefficient of \(t^{-1}dt\) in a Čech cocycle for the standard cover. Both steps commute with extension of scalars, hence so does the resulting pairing. Flat base change shows that \(g=\dim_kH^1(C,\mathcal O_C)\) is unchanged by field extension.
\end{proof}

% text: prose-C-07
\phantomsection\label{bp16}%
Recall that the chart of the Picard scheme consists of the invertible sheaves \(\mathcal L\) for which \(Rf_{T,*}\mathcal L\) is an invertible sheaf in degree \(0\). We use coherent base change in the following form.

% object: lemma-coherent-base-change-invertible-sheaves (blueprint/chap-The-Picard-scheme-of-a-curve/sec-Riemann-Roch-and-Serre-duality/lemma-coherent-base-change-invertible-sheaves/_section.md)
\phantomsection\label{bp17}%
\begin{bpstatementit}{Lemma~C.5}{ (Coherent base change for invertible sheaves)}
Let \(T\) be a \(k\)-scheme and let \(\mathcal L\) be invertible on \(C_T\). Then \(Rf_{T,*}\mathcal L\) is locally represented by a bounded complex of finite locally free \(\mathcal O_T\)-modules. For any \(a:T'\to T\), the base change map \(La^*Rf_{T,*}\mathcal L\to Rf_{T',*}\mathcal L_{T'}\) is an isomorphism in the derived category of \(\mathcal O_{T'}\)-modules. The locus where this complex is an invertible module in degree \(0\) is open. On this locus, \(f_{T,*}\mathcal L\) is invertible, its formation commutes with arbitrary base change, and \(R^1f_{T,*}\mathcal L=0\).
\end{bpstatementit}

% proof: proof/lemma-coherent-base-change-invertible-sheaves
\begin{proof}[{Proof}]\phantomsection\label{bp18}%
This is base change for coherent cohomology along a proper morphism, applied to a flat finitely presented module, together with the openness criterion for a perfect complex (see [Stacks, \href{https://stacks.math.columbia.edu/tag/0B91}{0B91}, \href{https://stacks.math.columbia.edu/tag/0B9S}{0B9S}]). To construct the complex, one may use a projective embedding, sufficiently large twists and a finite affine Čech complex; by Serre vanishing we then obtain models that are finite locally free. Given such a model, openness follows from a calculation with matrices. Indeed, near a point \(t\) where the derived fibre of the complex has cohomology only in degree \(0\), of dimension \(1\), we may cancel the differentials that are invertible at \(t\), and what remains is an invertible module in degree \(0\).

This proves the lemma for locally Noetherian \(T\). If \(T\) is an arbitrary \(k\)-scheme, then locally on \(T\) the sheaf \(\mathcal L\) is pulled back from \(C_{T_0}\) for some \(k\)-scheme \(T_0\) of finite type, and the base change isomorphism for \(T\to T_0\) reduces every assertion to \(T_0\). The open locus is the preimage of the locus for \(T_0\), and by base change it does not depend on \(T_0\).
\end{proof}


% object: sec-Hilbert-schemes-of-points (blueprint/chap-The-Picard-scheme-of-a-curve/sec-Hilbert-schemes-of-points/_section.md), content hidden
\subsection*{\texorpdfstring{C.3\quad Hilbert schemes of points}{C.3 Hilbert schemes of points}}
\phantomsection\label{bp19}%
\addcontentsline{toc}{subsection}{\texorpdfstring{C.3\quad Hilbert schemes of points}{C.3 Hilbert schemes of points}}
% text: prose-C-08
\phantomsection\label{bp20}%
On a smooth curve a finite subscheme of length \(d\) is the same as an effective divisor of degree \(d\). These divisors are parametrized by the Hilbert scheme of points \(\operatorname{Hilb}^{d}_{C/k}\), which is the symmetric power \(C^{(d)}\). Iterated addition of points gives a finite surjection \(C^d\to \operatorname{Hilb}^{d}_{C/k}\), from which \(\operatorname{Hilb}^{d}_{C/k}\) inherits properness, irreducibility and dimension \(d\) (Theorem \hyperref[bp21]{C.6}). For example, for \(C=\mathbf P^1\), an effective divisor of degree \(d\) is the zero set of a binary form of degree \(d\) up to scalars, and \(\operatorname{Hilb}^{d}_{C/k}=\mathbf P^d\).

A divisor \(D\) gives the invertible sheaf \(\mathcal O(D)\) with its canonical section. Conversely, an invertible sheaf with a section up to scalars gives a divisor. When \(h^0=1\) the section is unique up to scalars, and then the divisor is determined by the sheaf. In this way the open part of \(\operatorname{Hilb}^{g}_{C/k}\) where \(h^0(\mathcal O(D))=1\) and \(h^1=0\) becomes an open subfunctor of the Picard functor (Proposition \hyperref[bp24]{C.7}).

For \(d\geq0\), let \(\operatorname{Hilb}^d_{C/k}\) be the functor on \(k\)-schemes sending \(T\) to the set of closed subschemes \(D\subset C_T\) that are finite locally free of rank \(d\) over \(T\).

% object: theorem-hilbert-schemes-points-smooth-curve (blueprint/chap-The-Picard-scheme-of-a-curve/sec-Hilbert-schemes-of-points/theorem-hilbert-schemes-points-smooth-curve/_section.md)
\phantomsection\label{bp21}%
\begin{bpstatementit}{Theorem~C.6}{ (Hilbert schemes of points on a smooth curve)}
The functor \(\operatorname{Hilb}^{d}_{C/k}\) is represented by a smooth proper geometrically integral \(k\)-scheme of dimension \(d\), with \(\operatorname{Hilb}^{0}_{C/k}=\operatorname{Spec}k\). Its \(T\)-points are precisely the relative effective Cartier divisors that are finite locally free of degree \(d\). It has a universal divisor \(D_d\subset C\times_k\operatorname{Hilb}^{d}_{C/k}\) and an associated invertible sheaf \(\mathcal O(D_d)\).
\end{bpstatementit}

% proof: proof/theorem-hilbert-schemes-points-smooth-curve
\begin{proof}[{Proof}]\phantomsection\label{bp22}%
Representability is a general theorem, once we know that finite sets of points lie in affine opens. We then obtain the geometric properties from the addition map \(\operatorname{Hilb}^{d}_{C/k}\times C\to \operatorname{Hilb}^{d+1}_{C/k}\), which we show to be finite locally free. Starting from \(\operatorname{Hilb}^{1}_{C/k}=C\), we get smoothness and the dimension by induction on \(d\), and properness and irreducibility from the finite surjection \(C^d\to \operatorname{Hilb}^{d}_{C/k}\).

Every finite set of points of \(C\) lies in an affine open. Indeed, if \(Q\) is a closed point outside the set, then \(\mathcal O_C(Q)\) has positive degree, hence is ample by Riemann–Roch, and so \(C\setminus\{Q\}\) is affine. Hence \(\operatorname{Hilb}^{d}_{C/k}\) is represented by a scheme, separated and locally of finite type over \(k\) [Stacks, \href{https://stacks.math.columbia.edu/tag/0B9A}{0B9A}]. The universal closed subscheme is the divisor \(D_d\).

Let \(T\) be a \(k\)-scheme and \(t\in T\) a point. Since the local rings of \(C_{\kappa(t)}\) at closed points are discrete valuation rings, every finite closed subscheme of a fibre is an effective Cartier divisor. By the fibrewise criterion for relative effective Cartier divisors, every \(D\in \operatorname{Hilb}^{d}_{C/k}(T)\) is a relative effective Cartier divisor on \(C_T/T\). This gives the description of \(T\)-points and the invertible sheaf \(\mathcal O(D_d)\). Multiplying ideals adds such divisors, and we obtain the addition map

\[
\operatorname{Hilb}^{d}_{C/k}\times C\longrightarrow \operatorname{Hilb}^{d+1}_{C/k},\qquad (D,x)\longmapsto D+\Gamma_x.
\]

It identifies \(\operatorname{Hilb}^{d}_{C/k}\times C\) with the universal divisor over \(\operatorname{Hilb}^{d+1}_{C/k}\): if a divisor \(E\) contains the graph \(\Gamma_x\), then dividing a local equation of \(E\) by one of \(\Gamma_x\) gives a relative effective Cartier divisor \(D\) with \(E=D+\Gamma_x\), of degree \(d\) by the exact sequence for a sum of divisors. In other words, the addition map is the projection to \(\operatorname{Hilb}^{d+1}_{C/k}\) of the universal divisor \(D_{d+1}\subset C\times \operatorname{Hilb}^{d+1}_{C/k}\). By definition of the Hilbert functor, \(D_{d+1}\) is finite locally free of degree \(d+1\) over \(\operatorname{Hilb}^{d+1}_{C/k}\). Hence the addition map is finite locally free of degree \(d+1\) [Stacks, \href{https://stacks.math.columbia.edu/tag/0B9G}{0B9G}], in particular faithfully flat.

We now argue by induction on \(d\), starting from \(\operatorname{Hilb}^{0}_{C/k}=\operatorname{Spec}k\) and \(\operatorname{Hilb}^{1}_{C/k}=C\). The scheme \(\operatorname{Hilb}^{d+1}_{C/k}\) is smooth because regularity descends along faithfully flat maps, also after extending \(k\). It has dimension \(d+1\) because the addition map is finite and surjective [Stacks, \href{https://stacks.math.columbia.edu/tag/0B9H}{0B9H}].

Properness comes from the finite surjective map \(C^d\to \operatorname{Hilb}^{d}_{C/k}\) given by iterated addition. This map makes \(\operatorname{Hilb}^{d}_{C/k}\) quasi-compact. It also makes \(\operatorname{Hilb}^{d}_{C/k}\) universally closed, because \(C^d\) is proper and surjectivity is stable under base change. Since \(C^d\) is geometrically integral, \(\operatorname{Hilb}^{d}_{C/k}\) is geometrically irreducible. Being smooth, it is then geometrically integral [Stacks, \href{https://stacks.math.columbia.edu/tag/0B9I}{0B9I}].

For \(C=\mathbf P^1\), the map \(C^d\to \operatorname{Hilb}^{d}_{C/k}=\mathbf P^d\) sends \(d\) points to the binary form with these zeros, that is, the product of \(d\) linear forms; it is the quotient by the symmetric group, of degree \(d!\).
\end{proof}

% text: prose-C-09
\phantomsection\label{bp23}%
Assume now that \(C(k)\) contains a point \(P\), and let us describe \(\operatorname{Pic}_{C/k}\) by rigidified invertible sheaves as in Lemma \hyperref[bp7]{C.2}. We define a subfunctor \(F\subset\operatorname{Pic}_{C/k}\) as follows. A rigidified class \(\mathcal L\) on \(C_T\) lies in \(F(T)\) if, in the derived category of \(\mathcal O_T\)-modules,

\[
Rf_{T,*}\mathcal L\simeq\mathcal N[0]
\quad\text{for some invertible }\mathcal N\text{ on }T.
\]

By Lemma \hyperref[bp17]{C.5}, this condition defines an open locus of \(T\) compatible with base change, so that the inclusion \(F\subset\operatorname{Pic}_{C/k}\) is representable by open immersions.

% object: proposition-divisors-degree-g-open-chart (blueprint/chap-The-Picard-scheme-of-a-curve/sec-Hilbert-schemes-of-points/proposition-divisors-degree-g-open-chart/_section.md)
\phantomsection\label{bp24}%
\begin{bpstatementit}{Proposition~C.7}{ (Divisors of degree $g$ as an open chart)}
The functor \(F\) is represented by an open subscheme \(W\subset \operatorname{Hilb}^{g}_{C/k}\). The morphism \(W\to\operatorname{Pic}_{C/k}\) associates to a divisor its rigidified invertible sheaf. Consequently, \(F\) is smooth of dimension \(g\).
\end{bpstatementit}

% proof: proof/proposition-divisors-degree-g-open-chart
\begin{proof}[{Proof}]\phantomsection\label{bp25}%
We construct mutually inverse maps: a divisor gives its invertible sheaf, normalized along \(P\), and a family in \(F\) has, up to units on \(T\), a unique section in each fibre, whose zero scheme is a family of divisors of degree \(g\).

\emph{From divisors to sheaves.} On \(C\times \operatorname{Hilb}^{g}_{C/k}\), normalize the universal invertible sheaf as

\[
\mathcal L_g=\mathcal O(D_g)\otimes
f_{\operatorname{Hilb}^g}^*(P_{\operatorname{Hilb}^g}^*\mathcal O(D_g))^{-1}.
\]

Let \(W\) be the open locus where \(Rf_{\operatorname{Hilb}^g,*}\mathcal L_g\) is a line bundle in degree \(0\); then \(\mathcal L_g\) defines a morphism \(W\to F\).

\emph{From sheaves to divisors.} Let \(\mathcal L\) be a rigidified class in \(F(T)\) and set \(\mathcal N=f_{T,*}\mathcal L\). By assumption \(Rf_{T,*}\mathcal L\simeq\mathcal N'[0]\) for an invertible \(\mathcal N'\), so \(\mathcal N\simeq\mathcal N'\) is invertible. The evaluation map \(f_T^*\mathcal N\to\mathcal L\) is a section \(s\) of \(\mathcal L\otimes f_T^*\mathcal N^{-1}\). It is determined by \(\mathcal L\) up to a unit on \(T\), and such a unit does not change its zero scheme \(D\).

This \(D\) is a relative effective Cartier divisor. Indeed, the formation of \(\mathcal N\) commutes with base change (Lemma \hyperref[bp17]{C.5}). Hence on each fibre \(C_{\kappa(t)}\) the section \(s\) is a nonzero element of the one-dimensional space of sections. The fibre is integral because \(C\) is geometrically integral, so the zero scheme of \(s\) on the fibre is an effective Cartier divisor. Moreover, the fibres of \(D\) have degree \(g\): for a point \(t\) of \(T\) we have \(h^0(\mathcal L_t)=1\) and \(H^1(C_t,\mathcal L_t)=0\), so Riemann–Roch gives

\[
\deg\mathcal L_t
=\chi(\mathcal L_t)-(1-g)=1-(1-g)=g.
\]

The morphism \(D\to T\) is proper, since \(D\) is closed in the proper scheme \(C_T\). It has finite fibres, and it is flat, as is every relative effective Cartier divisor. A proper, flat, finitely presented morphism with finite fibres is finite locally free. Its rank at \(t\) is the degree \(g\) of the fibre \(D_t\), the divisor of zeros of a section of \(\mathcal L_t\). Hence \(D\) defines a morphism \(T\to \operatorname{Hilb}^{g}_{C/k}\).

\emph{The constructions are inverse.} Normalizing \(\mathcal O(D)\) along \(P_T\) recovers \(\mathcal L\), because \(\mathcal O(D)\simeq\mathcal L\otimes f_T^*\mathcal N^{-1}\). The two constructions are therefore inverse to each other, and they commute with base change, which proves that \(W\simeq F\) (see [Stacks, \href{https://stacks.math.columbia.edu/tag/0B9X}{0B9X}]).
\end{proof}

% text: prose-C-10
\phantomsection\label{bp26}%
To cover the Picard functor by translates of \(F\), we have to move every invertible sheaf into \(F\) by a twist defined over \(k\).

% object: lemma-twisting-h-0-1-h (blueprint/chap-The-Picard-scheme-of-a-curve/sec-Hilbert-schemes-of-points/lemma-twisting-h-0-1-h/_section.md)
\phantomsection\label{bp27}%
\begin{bpstatementit}{Lemma~C.8}{ (Twisting to $h^0=1$ and $h^1=0$)}
Suppose that \(k\) is separably closed. For every field extension \(k'/k\) and every invertible sheaf \(\mathcal L\) on \(C_{k'}\), there is an invertible sheaf \(\mathcal M_0\) on \(C\) such that

\[
h^0(C_{k'},\mathcal L\otimes\mathcal M_{0,k'})=1,
\qquad H^1(C_{k'},\mathcal L\otimes\mathcal M_{0,k'})=0.
\]
\end{bpstatementit}

% proof: proof/lemma-twisting-h-0-1-h
\begin{proof}[{Proof}]\phantomsection\label{bp28}%
We first twist by a large power of an ample sheaf on \(C\), which kills \(H^1\) and leaves many sections. Then we remove \(k\)-rational points one at a time, each time lowering \(h^0\) by one and keeping \(H^1=0\), until \(h^0=1\). The twist is defined over \(k\) because the points are.

Choose an ample invertible sheaf \(\mathcal M\) on \(C\). For \(r\gg0\), Serre vanishing and Riemann–Roch show that \(\mathcal L\otimes\mathcal M_{k'}^{\otimes r}\) has vanishing \(H^1\) and a positive-dimensional space of sections. As long as \(h^0\) is larger than \(1\), we choose a nonzero section and a point \(x\in C(k)\) where that section is nonzero. Such a point exists: the zero divisor of the section is finite, so its image in \(C\) is a finite set, and the complement of this set is a nonempty open of \(C\). This open contains a point of \(C(k)\), because the \(k\)-rational points of a smooth \(k\)-scheme are dense when \(k\) is separably closed [Stacks, \href{https://stacks.math.columbia.edu/tag/056U}{056U}]. For the current sheaf \(\mathcal F\), consider the sequence

\[
0\longrightarrow\mathcal F(-x_{k'})
\longrightarrow\mathcal F
\longrightarrow\mathcal F|_{x_{k'}}\longrightarrow0
\]

The chosen section does not vanish at \(x_{k'}\), so the map \(H^0(C_{k'},\mathcal F)\to\mathcal F|_{x_{k'}}\) is onto the one-dimensional fibre. By the long exact sequence, subtracting \(x_{k'}\) lowers \(h^0\) by exactly \(1\) and preserves \(H^1=0\). We continue in this way until \(h^0=1\). Since all subtracted points come from \(k\), the resulting twist has the form \(\mathcal M_0=\mathcal M^{\otimes r}(-x_1-\cdots-x_s)\) on \(C\). Finally, normalizing along \(P\) does not change the dimensions of fibre cohomology, so the rigidified twist satisfies the same conditions (see [Stacks, \href{https://stacks.math.columbia.edu/tag/0B9Y}{0B9Y}]).
\end{proof}


% object: sec-The-Jacobian (blueprint/chap-The-Picard-scheme-of-a-curve/sec-The-Jacobian/_section.md), content hidden
\subsection*{\texorpdfstring{C.4\quad The Jacobian}{C.4 The Jacobian}}
\phantomsection\label{bp29}%
\addcontentsline{toc}{subsection}{\texorpdfstring{C.4\quad The Jacobian}{C.4 The Jacobian}}
% text: prose-C-11
\phantomsection\label{bp30}%
We glue the translates of the chart of Proposition \hyperref[bp24]{C.7} to the Picard scheme, whose component of degree zero is the Jacobian. Gluing uses twice that \(k\) is separably closed: to find a rational point for the rigidification, and to obtain the twists of Lemma \hyperref[bp27]{C.8} over \(k\). For an elliptic curve, the Abel map \(C=\operatorname{Hilb}^{1}_{C/k}\to\operatorname{Pic}^1_{C/k}\) is an isomorphism, and translation by \(\mathcal O_C(-P)\) identifies \(C\) with its own Jacobian. In general, the tangent space of the Jacobian at the origin is computed by deforming the trivial sheaf over the dual numbers.

% object: lemma-representability-picard-functor (blueprint/chap-The-Picard-scheme-of-a-curve/sec-The-Jacobian/lemma-representability-picard-functor/_section.md)
\phantomsection\label{bp31}%
\begin{bpstatementit}{Lemma~C.9}{ (Representability of the Picard functor)}
Let \(k\) be separably closed and let \(C/k\) be a smooth proper connected curve. The fppf Picard functor \(\operatorname{Pic}_{C/k}\) is represented by a scheme locally of finite type over \(k\), which is covered by the open subschemes representing the translates \(M^{-1}F\), for \(M\in\operatorname{Pic}_{C/k}(k)\), of the chart \(F\) of Proposition \hyperref[bp24]{C.7}, defined using a point \(P\in C(k)\).
\end{bpstatementit}

% proof: proof/lemma-representability-picard-functor
\begin{proof}[{Proof}]\phantomsection\label{bp32}%
The preceding lemmas assume \(C\) geometrically integral. To see that this holds, note that \(\bar k/k\) is purely inseparable. Hence \(C_{\bar k}\to C\) is a homeomorphism, and \(C_{\bar k}\) is connected and smooth, hence integral. We rigidify using a point \(P\in C(k)\). Such a point exists because \(k\) is separably closed and the \(k\)-points of the smooth curve \(C\) are dense [Stacks, \href{https://stacks.math.columbia.edu/tag/056U}{056U}].

By Proposition \hyperref[bp24]{C.7} we have a representable open subfunctor \(F\subset\operatorname{Pic}_{C/k}\). For a fixed class \(M\in\operatorname{Pic}_{C/k}(k)\), the translate \(M^{-1}F\), consisting of the classes \(\mathcal L\) with \(\mathcal L\otimes M\) in \(F\), is another representable open subfunctor.

We show that these translates cover \(\operatorname{Pic}_{C/k}\). Let a class on a \(k\)-scheme \(T\) and a point \(t\in T\) be given. We apply Lemma \hyperref[bp27]{C.8} to its fibre over the residue field \(k'=\kappa(t)\). This gives an invertible sheaf \(\mathcal M_0\) on \(C\), defined over \(k\), such that the twisted fibre has \(h^0=1\) and \(H^1=0\). By Lemma \hyperref[bp17]{C.5}, the twisted class then lies in \(F\) over an open neighbourhood of \(t\). Letting \(t\) vary, we see that the open loci where the corresponding translated class lies in \(F\) cover \(T\). Note that the translates are indexed by classes over \(k\), while the points of \(T\) have arbitrary residue fields. This is why Lemma \hyperref[bp27]{C.8} asks for a twist defined over \(k\) for every extension \(k'\).

The intersections of the charts are open in either chart, and their transition isomorphisms satisfy the cocycle condition by functoriality. An fppf sheaf covered by representable open subfunctors is representable [Stacks, \href{https://stacks.math.columbia.edu/tag/0B9Q}{0B9Q}, \href{https://stacks.math.columbia.edu/tag/0B9Z}{0B9Z}]. Hence gluing the charts gives a scheme representing \(\operatorname{Pic}_{C/k}\). It is locally of finite type because each chart is an open of \(\operatorname{Hilb}^{g}_{C/k}\).
\end{proof}

% text: prose-C-12
\phantomsection\label{bp33}%
We compute the tangent space with the dual numbers \(k[\varepsilon]\), \(\varepsilon^2=0\); the computation works for any scheme over \(k\).

% object: lemma-tangent-space-picard-functor (blueprint/chap-The-Picard-scheme-of-a-curve/sec-The-Jacobian/lemma-tangent-space-picard-functor/_section.md)
\phantomsection\label{bp34}%
\begin{bpstatementit}{Lemma~C.10}{ (First-order deformations of invertible sheaves)}
Let \(X\) be a scheme over a field \(k\) and \(X_\varepsilon=X\times_k\operatorname{Spec}k[\varepsilon]\). The kernel of the restriction map \(\operatorname{Pic}(X_\varepsilon)\to\operatorname{Pic}(X)\) is \(H^1(X,\mathcal O_X)\).
\end{bpstatementit}

% proof: proof/lemma-tangent-space-picard-functor
\begin{proof}[{Proof}]\phantomsection\label{bp35}%
Because \(\varepsilon^2=0\), the map \(a\mapsto1+\varepsilon a\) is a homomorphism, and it gives an exact sequence of sheaves of groups on \(X\)

\[
0\longrightarrow\mathcal O_X\longrightarrow
\mathcal O_{X_\varepsilon}^\times\longrightarrow\mathcal O_X^\times
\longrightarrow1,
\]

split by the inclusion \(\mathcal O_X^\times\subset
\mathcal O_{X_\varepsilon}^\times\). The schemes \(X\) and \(X_\varepsilon\) have the same underlying space, so \(\operatorname{Pic}(X_\varepsilon)=H^1(X,\mathcal O_{X_\varepsilon}^\times)\). Since the sequence is split, taking \(H^1\) gives \(H^1(X,\mathcal O_{X_\varepsilon}^\times)=H^1(X,\mathcal O_X)\oplus\operatorname{Pic}(X)\), in which the restriction map is the projection to the second summand. Hence the kernel is \(H^1(X,\mathcal O_X)\).
\end{proof}

% object: theorem-picard-scheme-jacobian-over-separably (blueprint/chap-The-Picard-scheme-of-a-curve/sec-The-Jacobian/theorem-picard-scheme-jacobian-over-separably/_section.md)
\phantomsection\label{bp36}%
\begin{bpstatementit}{Theorem~C.11}{ (The Picard scheme and the Jacobian over a separably closed field)}
Let \(k\) be separably closed and let \(C/k\) be a smooth proper connected curve of genus \(g\). The fppf Picard functor is represented by a smooth separated group scheme, locally of finite type over \(k\), with a decomposition

\[
\operatorname{Pic}_{C/k}
=\coprod_{d\in\mathbf Z}\operatorname{Pic}^d_{C/k}.
\]

Each component is a smooth proper geometrically connected variety of dimension \(g\), parametrizing invertible sheaves of degree \(d\). The identity component \(J(C)=\operatorname{Pic}^0_{C/k}\) is an abelian variety, and

\[
J(C)(k)=\operatorname{Pic}^0(C),
\qquad T_0J(C)=H^1(C,\mathcal O_C).
\]

The formation of this representing scheme, its degree components and \(J(C)\) commutes with extension of the ground field.
\end{bpstatementit}

% proof: proof/theorem-picard-scheme-jacobian-over-separably
\begin{proof}[{Proof}]\phantomsection\label{bp37}%
Representability is Lemma \hyperref[bp31]{C.9}. The tangent space is given by Lemma \hyperref[bp34]{C.10}. Fix the point \(P\in C(k)\) used to define \(F\).

For an invertible sheaf \(\mathcal L\) on \(C_T\), the Euler characteristic \(\chi(C_t,\mathcal L_t)\) is a locally constant function of \(t\in T\). This holds because \(C_T\to T\) is proper and flat and \(Rf_{T,*}\mathcal L\) is perfect (Lemma \hyperref[bp17]{C.5}). By Riemann–Roch the degree of \(\mathcal L_t\) is locally constant as well. This gives the displayed decomposition into open and closed pieces. The charts in \(\operatorname{Hilb}^{g}_{C/k}\) and their translates show that the scheme is smooth of dimension \(g\).

It remains to see that the scheme is separated. A \(k\)-rational point of a scheme locally of finite type over a field is closed. Hence the unit section of the group scheme is a closed immersion. The diagonal is the preimage of the unit section under \((x,y)\mapsto xy^{-1}\), so it is closed as well.

The universal divisor defines the Abel map

\[
\operatorname{Hilb}^{g}_{C/k}\longrightarrow\operatorname{Pic}^g_{C/k},
\qquad D\longmapsto\mathcal O_C(D).
\]

For a line bundle \(\mathcal L\) of degree \(g\) on \(C_{k'}\), with \(k'\) algebraically closed, Riemann–Roch gives \(h^0(\mathcal L)\geq\chi(\mathcal L)=g+1-g=1\). A nonzero section has an effective divisor \(D\) of degree \(g\) with \(\mathcal O(D)\simeq\mathcal L\). Hence the Abel map is surjective on points over algebraically closed fields, and this suffices for surjectivity. It follows that \(\operatorname{Pic}^g_{C/k}\) is quasi-compact. It is of finite type since the charts are locally of finite type. A separated scheme of finite type that is the surjective image of the proper scheme \(\operatorname{Hilb}^{g}_{C/k}\) is proper. It is geometrically irreducible because \(\operatorname{Hilb}^{g}_{C/k}\) is. Translation by \(\mathcal O_C((d-g)P)\) identifies \(\operatorname{Pic}^g_{C/k}\) with the component of each degree \(d\). In particular, the component of degree zero is smooth, proper, connected and commutative. Such a group scheme is an abelian variety (see [Stacks, \href{https://stacks.math.columbia.edu/tag/0BA0}{0BA0}]).

The tangent space \(T_0J(C)\) is the kernel of \(\operatorname{Pic}_{C/k}(k[\varepsilon])\to\operatorname{Pic}_{C/k}(k)\), where \(\varepsilon^2=0\). Since \(\operatorname{Pic}(k[\varepsilon])=0\), this is the kernel of \(\operatorname{Pic}(C_{k[\varepsilon]})\to\operatorname{Pic}(C)\), which is \(H^1(C,\mathcal O_C)\) by Lemma \hyperref[bp34]{C.10}.

For a field extension \(k'/k\), the Picard functor of \(C_{k'}\) is the restriction of \(\operatorname{Pic}_{C/k}\) to \(k'\)-schemes. Hence its representing scheme is the base change. The degree of \(\mathcal L_t\) is computed on the fibre \(C_t\), so the decomposition by degree, and with it \(J(C)\), is compatible with this base change.
\end{proof}

% text: prose-C-13
\phantomsection\label{bp38}%
We use the theorem in the following form: over an algebraically closed field, the Picard group of \(C\) is an extension of \(\mathbf Z\) by the group of points of an abelian variety.

% object: corollary-degree-jacobian (blueprint/chap-The-Picard-scheme-of-a-curve/sec-The-Jacobian/corollary-degree-jacobian/_section.md)
\phantomsection\label{bp39}%
\begin{bpstatementit}{Corollary~C.12}{ (Degree and the Jacobian)}
Let \(C\) be a smooth proper connected curve over an algebraically closed field \(k\). There is an exact sequence

\[
0\longrightarrow J(C)(k)
\longrightarrow\operatorname{Pic}(C)
\xrightarrow{\deg}\mathbf Z\longrightarrow0,
\]

and a point \(P\in C(k)\) splits it by \(d\mapsto\mathcal O_C(dP)\).
\end{bpstatementit}

% proof: proof/corollary-degree-jacobian
\begin{proof}[{Proof}]\phantomsection\label{bp40}%
The degree is additive under tensor product, so it is a homomorphism. By Theorem \hyperref[bp36]{C.11}, its kernel is \(\operatorname{Pic}^0(C)=J(C)(k)\). Since \(\deg\mathcal O_C(dP)=d\), the map \(d\mapsto\mathcal O_C(dP)\) is a section of the degree, which is therefore surjective.
\end{proof}

% object: example-degree-sequence-genus-zero-one (blueprint/chap-The-Picard-scheme-of-a-curve/sec-The-Jacobian/example-degree-sequence-genus-zero-one/_section.md)
\phantomsection\label{bp41}%
\begin{bpstatement}{Example~C.13}{ (The degree sequence in genus zero and one)}
On \(C=\mathbf P^1_k\) a divisor \(\sum_in_ia_i+n_\infty\infty\) of degree zero, with \(a_i\in k\), is the divisor of the rational function \(\prod_i(x-a_i)^{n_i}\), whose order at \(\infty\) is \(-\sum_in_i=n_\infty\). So \(\operatorname{Pic}^0(\mathbf P^1)=0\), and the degree is an isomorphism \(\operatorname{Pic}(\mathbf P^1)\simeq\mathbf Z\). This agrees with \(J(\mathbf P^1)=0\) for \(g=0\).

On an elliptic curve \(E\) with origin \(O\) the kernel is as large as the curve. Let \(\mathcal L\) have degree zero. Then \(\mathcal L(O)\) has degree one, and by Riemann–Roch with \(g=1\) it has a one-dimensional space of sections. Hence \(\mathcal L(O)\simeq\mathcal O_E(P)\) for a unique point \(P\), and \(P\mapsto\mathcal O_E(P-O)\) is a bijection \(E(k)\to\operatorname{Pic}^0(E)\). In particular \(\mathcal O_E(P-Q)\) is nontrivial for \(P\neq Q\). One can also see this directly: a rational function with divisor \(P-Q\) would be a map \(E\to\mathbf P^1\) of degree one, that is, an isomorphism.
\end{bpstatement}


% object: sec-Torsion-points (blueprint/chap-The-Picard-scheme-of-a-curve/sec-Torsion-points/_section.md), content hidden
\subsection*{\texorpdfstring{C.5\quad Torsion points}{C.5 Torsion points}}
\phantomsection\label{bp42}%
\addcontentsline{toc}{subsection}{\texorpdfstring{C.5\quad Torsion points}{C.5 Torsion points}}
% text: prose-C-14
\phantomsection\label{bp43}%
The Kummer sequence on a curve \(C\) over an algebraically closed field identifies \(H^1(C,\mu_n)\) with \(\operatorname{Pic}(C)[n]\) and \(H^2(C,\mu_n)\) with \(\operatorname{Pic}(C)/n\) (Proposition \bpout{D.2}). To compute these groups, we combine two facts. First, the exact sequence of Corollary \hyperref[bp39]{C.12} separates the discrete part \(\mathbf Z\) of the Picard group from the Jacobian. Second, Appendix \bpout{B} describes multiplication by \(n\) on an abelian variety over an algebraically closed field: it is surjective on points, and for \(n\) invertible in \(k\) its kernel is \((\mathbf Z/n)^{2g}\). Now apply the snake lemma to multiplication by \(n\) on the degree sequence. Since \(\mathbf Z\) is torsion free, the \(n\)-torsion of \(\operatorname{Pic}(C)\) is \(J(C)[n](k)\simeq(\mathbf Z/n)^{2g}\). Since \(J(C)(k)\) is divisible, the cokernel of \(n\) on \(\operatorname{Pic}(C)\) is \(\mathbf Z/n\), given by the degree.

For proper base change we also need these groups to be unchanged under extension of the algebraically closed field. This holds because \(J(C)[n]\) is finite étale (Corollary \hyperref[bp45]{C.15}).

% object: example-two-torsion-elliptic-curve (blueprint/chap-The-Picard-scheme-of-a-curve/sec-Torsion-points/example-two-torsion-elliptic-curve/_section.md)
\phantomsection\label{bp44}%
\begin{bpstatement}{Example~C.14}{ (Two-torsion of an elliptic curve)}
Let \(k\) be algebraically closed of characteristic different from two and let \(E\colon y^2=(x-e_1)(x-e_2)(x-e_3)\) with distinct \(e_i\in k\), with origin \(O\) at infinity and \(P_i=(e_i,0)\). The function \(x-e_i\) has a double zero at \(P_i\), because \(y\) is a local parameter there. It has a double pole at \(O\). Hence \(\operatorname{div}(x-e_i)=2P_i-2O\), and the classes \(\mathcal O_E(P_i-O)\) are killed by \(2\). They are nonzero and pairwise distinct, since \(\mathcal O_E(P_i-P_j)\) is trivial only for \(i=j\) (Example \hyperref[bp41]{C.13}). The divisor \(\operatorname{div}(y)=P_1+P_2+P_3-3O\) shows that the three classes add up to zero. So \(0\) and the classes of \(P_i-O\) form a subgroup \(\{0,P_1-O,P_2-O,P_3-O\}\simeq(\mathbf Z/2)^2\) of \(\operatorname{Pic}(E)[2]\). By Corollary \hyperref[bp45]{C.15} below, \(\operatorname{Pic}(E)[2]\) has order \(2^{2g}=4\), so this subgroup is all of it.

Under the Kummer sequence, the class \(P_i-O\), with the trivialization of its square given by \(x-e_i\), is the \(\mu_2\)-torsor obtained by adjoining a square root of \(x-e_i\). This torsor is unramified everywhere because the divisor of \(x-e_i\) is even. For the running example \(y^2=x^3-x\) over \(\overline{\mathbf F}_3\) the \(e_i\) are \(0,1,-1\in\mathbf F_3\). Hence the three torsors are defined over \(\mathbf F_3\), and Frobenius fixes all of \(E[2]\), as noted in Section \bpout{B.3}.
\end{bpstatement}

% object: corollary-torsion-jacobian-field-extension (blueprint/chap-The-Picard-scheme-of-a-curve/sec-Torsion-points/corollary-torsion-jacobian-field-extension/_section.md)
\phantomsection\label{bp45}%
\begin{bpstatementit}{Corollary~C.15}{ (Torsion of the Jacobian and field extension)}
Let \(k\) be algebraically closed, let \(C/k\) be a smooth proper connected curve of genus \(g\), and let \(n>0\) be invertible in \(k\). Then there is a noncanonical isomorphism

\[
J(C)[n](k)\simeq(\mathbf Z/n\mathbf Z)^{2g}.
\]

For every algebraically closed extension \(k'/k\), the natural map

\[
J(C)[n](k)\longrightarrow J(C_{k'})[n](k')
\]

is an isomorphism. Multiplication by \(n\) on \(J(C)(k)\) is surjective, and consequently

\[
\operatorname{Pic}(C)[n]=J(C)[n](k),
\qquad
\operatorname{Pic}(C)/n\operatorname{Pic}(C)
\xrightarrow[\deg]{\ \sim\ }\mathbf Z/n\mathbf Z.
\]
\end{bpstatementit}

% proof: proof/corollary-torsion-jacobian-field-extension
\begin{proof}[{Proof}]\phantomsection\label{bp46}%
The assertions about \(J(C)\) are those of Appendix \bpout{B} for this abelian variety, and the degree sequence carries them over to \(\operatorname{Pic}(C)\).

By Proposition \bpout{B.12}, the kernel \(J(C)[n]\) is finite étale of rank \(n^{2g}\). A finite étale scheme of rank \(r\) over the algebraically closed field \(k\) is a disjoint union of \(r\) copies of \(\operatorname{Spec}k\), a description stable under base extension. By Theorem \hyperref[bp36]{C.11} the formation of the Jacobian and of \([n]\) commutes with base change, so the natural map \(J(C)[n](k)\to J(C_{k'})[n](k')\) is a bijection.

If \(\ell^m\) is the exact power of a prime \(\ell\) dividing \(n\), then \(J(C)[\ell^m](k)\simeq(\mathbf Z/\ell^m)^{2g}\) by Proposition \bpout{B.14} applied to \(J(C)\). Combining the parts for the primes dividing \(n\), we get the stated \(\mathbf Z/n\mathbf Z\)-module structure.

By Proposition \bpout{B.12}, the group \(J(C)(k)\) is divisible. In the exact sequence of Corollary \hyperref[bp39]{C.12}, the group \(\mathbf Z\) has no \(n\)-torsion and \(J(C)(k)/n=0\). This proves the two assertions for \(\operatorname{Pic}(C)\).
\end{proof}


% object: sec-Descent-to-the-ground-field (blueprint/chap-The-Picard-scheme-of-a-curve/sec-Descent-to-the-ground-field/_section.md), content hidden
\subsection*{\texorpdfstring{C.6\quad Descent to the ground field}{C.6 Descent to the ground field}}
\phantomsection\label{bp47}%
\addcontentsline{toc}{subsection}{\texorpdfstring{C.6\quad Descent to the ground field}{C.6 Descent to the ground field}}
% text: prose-C-15
\phantomsection\label{bp48}%
To compare the Galois action on the cohomology of a curve over \(\bar k\) with a Tate module, we need the Jacobian over the ground field \(k\). For example, in Proposition \bpout{14.3} we work over a finite field, and Frobenius acts on \(\varprojlim_nH^1(C,\mathbf Z/\ell^n(1))\). It has no fixed vector there because this module is \(T_\ell J_0\) for an abelian variety \(J_0\) over \(\mathbf F_q\) (Proposition \bpout{B.15}). The construction so far works over a separably closed field and uses a rational point for the rigidification. However, a curve over \(k\) need not have a \(k\)-point; for instance, the conic \(x^2+y^2+z^2=0\) over \(\mathbf R\) has none. We therefore define \(J_0\) as the scheme representing the degree-zero part of the fppf Picard sheaf. Concretely, we obtain it from the Jacobian over \(k_s\) by Galois descent through a finite Galois extension \(k'/k\) (Proposition \hyperref[bp49]{C.16}). The Kummer sequence then identifies \(H^1(C,\mathbf Z/\ell^n(1))\) with \(J_0[\ell^n](\bar k)\) compatibly with the Galois action and the transition maps (Proposition \hyperref[bp52]{C.17}).

% object: proposition-jacobian-over-ground-field (blueprint/chap-The-Picard-scheme-of-a-curve/sec-Descent-to-the-ground-field/proposition-jacobian-over-ground-field/_section.md)
\phantomsection\label{bp49}%
\begin{bpstatementit}{Proposition~C.16}{ (The Jacobian over the ground field)}
Let \(C_0/k\) be a smooth proper geometrically connected curve of genus \(g\), and put \(C=C_0\times_k\bar k\). There is an abelian variety \(J_0/k\) of dimension \(g\) representing the degree-zero relative Picard fppf sheaf of \(C_0\). For every field extension \(k'/k\), the base change \(J_0\times_kk'\) represents the degree-zero Picard sheaf of \(C_0\times_kk'\); in particular \(J_0\times_k\bar k\) represents \(\operatorname{Pic}^0_{C/\bar k}\), and the canonical identification

\[
J_0(\bar k)=\operatorname{Pic}^0(C)
\]

is Galois-equivariant.
\end{bpstatementit}

% proof: proof/proposition-jacobian-over-ground-field
\begin{proof}[{Proof}]\phantomsection\label{bp50}%
We descend the Jacobian of \(C_{0,k_s}\), where \(k_s\subseteq\bar
k\) is a separable closure of \(k\). In general \(k_s/k\) is infinite, hence not an fppf cover. For this reason we move the descent datum on the Jacobian over \(k_s\) by limit arguments to a finite Galois extension, and descend along that extension. A second limit argument shows that the descended scheme represents the degree-zero Picard sheaf.

Let \(F=\operatorname{Pic}^0_{C_0/k}\) be the degree-zero part of the fppf Picard sheaf. Here we use that the fibrewise degree is locally constant, by coherent base change and Riemann–Roch. Since the Picard sheaf and the degree commute with base change, Theorem \hyperref[bp36]{C.11} gives an abelian variety \(J_s\) over \(k_s\) representing the restriction of \(F\) to \(k_s\)-schemes; \(J_s\) is projective. Over \(k_s\otimes_kk_s\) the two pullbacks of \(J_s\) represent the same functor, so by the Yoneda lemma we get an isomorphism between them. All identifications are induced by \(F\). Hence this isomorphism satisfies the cocycle condition over \(k_s\otimes_kk_s\otimes_kk_s\) and respects the group structures.

To move this descent datum to a finite extension, note that \(k_s\) is the filtered union of its finite separable subextensions, and \(k_s\otimes_kk_s\) and the triple tensor product are the corresponding unions. By the limit theorems for schemes of finite presentation [Stacks, \href{https://stacks.math.columbia.edu/tag/01ZM}{01ZM}], the defining equations of \(J_s\) in a projective space, the descent isomorphism, its inverse and the group operations are defined over some finite stage. The finitely many identities among them (inverse, cocycle, group axioms) hold after enlarging the stage. Finally, we enlarge it further to a Galois extension. In this way we obtain a finite Galois extension \(K/k\) in \(k_s\) and a projective \(K\)-scheme \(J_K\) with a group structure and a compatible descent datum for \(K/k\), whose base change to \(k_s\) is the datum on \(J_s\). Since \(\operatorname{Spec}K\to\operatorname{Spec}k\) is finite locally free and surjective, the descent datum on the projective scheme \(J_K\) is effective [Stacks, \href{https://stacks.math.columbia.edu/tag/0CCJ}{0CCJ}]. This gives a commutative group scheme \(J_0\) over \(k\) with \(J_0\times_kk_s\simeq J_s\). Properness and smoothness descend along \(K/k\), and connectedness and dimension may be checked over \(\bar
k\). Hence \(J_0\) is an abelian variety of dimension \(g\).

The functors \(J_0\) and \(F\) are fppf sheaves that agree after base change to \(k_s\). Since \(k_s/k\) is not finite, we need a limit argument to show that they agree. Both functors are locally of finite presentation. For \(J_0\) this is [Stacks, \href{https://stacks.math.columbia.edu/tag/01ZC}{01ZC}]. For \(F\) it follows from Lemma \hyperref[bp7]{C.2} once the finite stages contain a field of definition of a point of \(C_0\), because invertible sheaves, rigidifications and their isomorphisms are of finite presentation. Hence, for affine \(T\) over \(k\), a section of either functor over \(T_{k_s}\), or over \(T_{k_s\otimes_kk_s}\), comes from a finite stage. Likewise, two sections that agree over \(T_{k_s}\) agree over a finite stage.

The identification \(J_0(T_{k_s})\simeq F(T_{k_s})\) is compatible with the descent data. Let a section of one functor over \(T\) be given. We transport its image over \(T_{k_s}\) and descend the result to some \(T_L\) with \(L/k\) finite. After enlarging \(L\), its two pullbacks to \(T_{L\otimes_kL}\) agree. Since both functors are fppf sheaves, it then descends to \(T\). The same argument applied to equalities shows that the result is unique. We obtain inverse bijections \(J_0(T)\simeq F(T)\), natural in affine \(T\), and hence an isomorphism of sheaves.

The Picard sheaf and the degree commute with field extension. Hence \(J_0\times_kk'\) represents the degree-zero Picard sheaf of \(C_0\times_kk'\) for every \(k'/k\). Over \(\bar k\) the value of the Picard sheaf is the Picard group, by rigidification and \(\operatorname{Pic}(\bar k)=0\). This gives the displayed identification. For \(\gamma\in\operatorname{Gal}(\bar k/k)\), naturality with respect to \(\operatorname{Spec}(\gamma)\) says that this group isomorphism intertwines the Galois action on points with pullback of invertible sheaves. Here \(\operatorname{Gal}(\bar k/k)=\operatorname{Gal}(k_s/k)\) because \(\bar k/k_s\) is purely inseparable.
\end{proof}

% text: prose-C-16
\phantomsection\label{bp51}%
Since \(J_0\) is defined over \(k\), the Galois group acts on its torsion points. Using the Kummer sequence on \(C\), we compare this action with the Galois action on cohomology.

% object: proposition-h-1-curve-tate-module (blueprint/chap-The-Picard-scheme-of-a-curve/sec-Descent-to-the-ground-field/proposition-h-1-curve-tate-module/_section.md)
\phantomsection\label{bp52}%
\begin{bpstatementit}{Proposition~C.17}{ ($H^1$ of a curve and the Tate module of its Jacobian)}
In the preceding setting, suppose \(\ell\) is invertible in \(k\). For every \(n\geq1\) there is a canonical Galois-equivariant isomorphism, with cohomology computed on the pro-étale site,

\[
H^1(C,\mathbf Z/\ell^n(1))
\simeq J_0[\ell^n](\bar k).
\]

The coefficient map \(\mu_{\ell^{n+1}}\to\mu_{\ell^n}\), \(z\mapsto z^\ell\), corresponds to multiplication by \(\ell\) on the torsion points. Consequently

\[
\varprojlim_n H^1(C,\mathbf Z/\ell^n(1))\simeq T_\ell J_0
\]

as \(\mathbf Z_\ell\)-modules with Galois action.
\end{bpstatementit}

% proof: proof/proposition-h-1-curve-tate-module
\begin{proof}[{Proof}]\phantomsection\label{bp53}%
The isomorphism comes from the Kummer sequence and the identification of the preceding proposition. The point is that both are compatible with the Galois action and with the transition maps. Then they pass to the inverse limit.

We compute on the étale site and transport the result to the pro-étale site by Theorem \bpout{2.35}. Since that comparison is natural, it preserves all the maps and Galois actions below. The Kummer sequence gives

\[
0\longrightarrow
H^0(C,\mathbf G_m)/\ell^n
\longrightarrow H^1(C,\mu_{\ell^n})
\longrightarrow\operatorname{Pic}(C)[\ell^n]
\longrightarrow0.
\]

The first term vanishes. Indeed, \(C\) is proper and connected, so \(H^0(C,\mathbf G_m)=\bar k^{\times}\), and this group is divisible. For the last term, note that the degree map takes values in the torsion-free group \(\mathbf Z\), so \(\operatorname{Pic}(C)[\ell^n]=\operatorname{Pic}^0(C)[\ell^n]\), and Proposition \hyperref[bp49]{C.16} identifies this group with \(J_0[\ell^n](\bar k)\). It remains to check Galois equivariance. The map induced by \(\mu_{\ell^n}\to\mathbf G_m\) and the identification of Proposition \hyperref[bp49]{C.16} are natural in the curve and in the field, and hence Galois-equivariant.

To check the transition maps, we map the Kummer sequence for \(\ell^{n+1}\) to that for \(\ell^n\) by \([\ell]\) on the first \(\mathbf G_m\) and by the identity on the second. On the kernels this is \(z\mapsto z^\ell\). On the Picard group it sends an invertible sheaf to its \(\ell\)-th tensor power, which is \([\ell]\) on \(J_0\). Therefore the displayed isomorphisms identify the two inverse systems, including their transition maps.
\end{proof}



\clearpage\phantomsection\hypertarget{bp-outside}{}%
\section*{Outside this export}
\addcontentsline{toc}{section}{Outside this export}
References marked \ensuremath{{}^\ast} are to parts of the blueprint this document leaves out:
\begin{description}
\item[Theorem 2.35] Bounded-below cohomology comparison, in Ⅰ The six operations (\bpslug{theorem-{}\allowbreak{}bounded-{}\allowbreak{}below-{}\allowbreak{}cohomology-{}\allowbreak{}comparison})
\item[Proposition 14.3] No fixed class in degree one, in Ⅲ Weights and the Riemann hypothesis (\bpslug{proposition-{}\allowbreak{}no-{}\allowbreak{}fixed-{}\allowbreak{}class-{}\allowbreak{}degree-{}\allowbreak{}one})
\item[Appendix B] Abelian varieties, in Appendices (\bpslug{chap-{}\allowbreak{}Abelian-{}\allowbreak{}varieties})
\item[Proposition B.12] Multiplication by $n$, in Appendices (\bpslug{proposition-{}\allowbreak{}multiplication-{}\allowbreak{}n})
\item[Section B.3] Tate modules, in Appendices (\bpslug{sec-{}\allowbreak{}Tate-{}\allowbreak{}modules})
\item[Proposition B.14] The integral Tate module, in Appendices (\bpslug{proposition-{}\allowbreak{}integral-{}\allowbreak{}tate-{}\allowbreak{}module})
\item[Proposition B.15] No Frobenius-fixed vector, in Appendices (\bpslug{proposition-{}\allowbreak{}no-{}\allowbreak{}frobenius-{}\allowbreak{}fixed-{}\allowbreak{}vector})
\item[Proposition D.2] Smooth proper curves, torsion prime to the characteristic, in Appendices (\bpslug{proposition-{}\allowbreak{}smooth-{}\allowbreak{}proper-{}\allowbreak{}curves-{}\allowbreak{}torsion-{}\allowbreak{}prime})
\end{description}
\clearpage\phantomsection%
\section*{References}
\addcontentsline{toc}{section}{References}
\bpbibitem{BBD82}{bp55}{A. Beilinson, J. Bernstein, and P. Deligne. \emph{Faisceaux pervers}, volume 100 of \emph{Astérisque}. 1982.}

\bpbibitem{BS15}{bp56}{B. Bhatt and P. Scholze. The pro-étale topology for schemes. \emph{Astérisque}, 369:99–201, 2015. \href{https://people.mpim-bonn.mpg.de/scholze/proetale.pdf}{Author's manuscript}.}

\bpbibitem{Con07}{bp57}{B. Conrad. Deligne's notes on Nagata compactifications. \emph{Journal of the Ramanujan Mathematical Society}, 22:205–257, 2007. \href{https://math.stanford.edu/\~conrad/papers/nagatafinal.pdf}{Author's manuscript}.}

\bpbibitem{Del74}{bp58}{P. Deligne. La conjecture de Weil. I. \emph{Publ. Math. IHÉS}, 43:273–307, 1974.}

\bpbibitem{Del80}{bp59}{P. Deligne. La conjecture de Weil. II. \emph{Publ. Math. IHÉS}, 52:137–252, 1980.}

\bpbibitem{Den22}{bp60}{C. Deninger. There is no “Weil-”cohomology theory with real coefficients for arithmetic curves. \emph{Ann. Sc. Norm. Super. Pisa Cl. Sci.}, 2022. To appear. arXiv:\href{https://arxiv.org/abs/2204.02714}{2204.02714}.}

\bpbibitem{Deu41}{bp61}{M. Deuring. Die Typen der Multiplikatorenringe elliptischer Funktionenkörper. \emph{Abh. Math. Sem. Univ. Hamburg}, 14:197–272, 1941.}

\bpbibitem{Eke90}{bp62}{T. Ekedahl. On the adic formalism. In \emph{The Grothendieck Festschrift, Vol. II}, volume 87 of \emph{Progress in Mathematics}, pages 197–218. Birkhäuser, 1990.}

\bpbibitem{Has36}{bp63}{H. Hasse. Zur Theorie der abstrakten elliptischen Funktionenkörper. I, II, III. \emph{J. Reine Angew. Math.}, 175:55–62, 69–88, 193–208, 1936.}

\bpbibitem{Hum72}{bp64}{J. E. Humphreys. \emph{Introduction to Lie Algebras and Representation Theory}, volume 9 of \emph{Graduate Texts in Mathematics}. Springer, 1972.}

\bpbibitem{Ill94}{bp65}{L. Illusie. Autour du théorème de monodromie locale. \emph{Astérisque}, 223:9–57, 1994. §4. \href{https://www.numdam.org/item/AST_1994__223__9_0/}{Original article}.}

\bpbibitem{Kat88}{bp66}{N. M. Katz. Travaux de Laumon. \emph{Astérisque}, 161–162:105–132, 1988. Séminaire Bourbaki, exp. 691.}

\bpbibitem{KW01}{bp67}{R. Kiehl and R. Weissauer. \emph{Weil conjectures, perverse sheaves and \(\ell\)-adic Fourier transform}, volume 42 of \emph{Ergebnisse der Mathematik}. Springer, 2001.}

\bpbibitem{Lau87}{bp68}{G. Laumon. Transformation de Fourier, constantes d'équations fonctionnelles et conjecture de Weil. \emph{Publ. Math. IHÉS}, 65:131–210, 1987.}

\bpbibitem{Mil86a}{bp69}{J. S. Milne. \href{https://www.jmilne.org/math/xnotes/AVs.pdf}{Abelian Varieties}. In \emph{Arithmetic Geometry}. 1986. Corrected online version.}

\bpbibitem{Mil86b}{bp54}{J. S. Milne. \href{https://www.jmilne.org/math/xnotes/JVs.pdf}{Jacobian Varieties}. In \emph{Arithmetic Geometry}. 1986. Corrected online version.}

\bpbibitem{Ran39}{bp70}{R. A. Rankin. Contributions to the theory of Ramanujan's function \(\tau(n)\) and similar arithmetical functions. II. \emph{Proc. Cambridge Philos. Soc.}, 35:357–372, 1939.}

\bpbibitem{SGA4}{bp71}{M. Artin, A. Grothendieck, and J.-L. Verdier, editors. \emph{Théorie des topos et cohomologie étale des schémas}, volumes 269, 270, 305 of \emph{Lecture Notes in Mathematics}. Springer, 1972–1973. SGA 4.}

\bpbibitem{SGA4½}{bp72}{P. Deligne. \emph{Cohomologie étale}, volume 569 of \emph{Lecture Notes in Mathematics}. Springer, 1977. SGA 4½; \emph{Théorèmes de finitude}, §3, and \emph{Rapport sur la formule des traces}, §3, and \emph{Sommes trigonométriques}, §1.}

\bpbibitem{SGA5}{bp73}{A. Grothendieck, editor. \emph{Cohomologie \(\ell\)-adique et fonctions \(L\)}, volume 589 of \emph{Lecture Notes in Mathematics}. Springer, 1977. SGA 5.}

\bpbibitem{SGA7II}{bp74}{P. Deligne and N. Katz, editors. \emph{Groupes de monodromie en géométrie algébrique}, volume 340 of \emph{Lecture Notes in Mathematics}. Springer, 1973. SGA 7 II; Exposé XIII, (2.1.7).}

\bpbibitem{Stacks}{bp75}{The Stacks Project Authors. \href{https://stacks.math.columbia.edu}{The Stacks Project}. Individual results are cited by their stable tags in the text.}

\bpbibitem{Var05}{bp76}{Y. Varshavsky. Lefschetz–Verdier trace formula and a generalization of a theorem of Fujiwara. 2005. Version 2. arXiv:\href{https://arxiv.org/abs/math/0505564}{math/0505564}. \href{https://arxiv.org/abs/math/0505564v2}{Author's preprint}.}

\bpbibitem{Wei48a}{bp77}{A. Weil. \emph{Sur les courbes algébriques et les variétés qui s'en déduisent}. Hermann, Paris, 1948.}

\bpbibitem{Wei48b}{bp78}{A. Weil. \emph{Variétés abéliennes et courbes algébriques}. Hermann, Paris, 1948.}

\bpbibitem{Wei49}{bp79}{A. Weil. Numbers of solutions of equations in finite fields. \emph{Bull. Amer. Math. Soc.}, 55:497–508, 1949.}

\end{document}
